\documentclass[11pt,a4paper]{article}
\usepackage[margin=24mm]{geometry}
\usepackage{amsmath,amssymb}
\title{Class Field Theory VII.9.3}
\author{J. S. Milne}
\date{Fixed author PDF v4.03}
\begin{document}
\maketitle
\section*{Original statement}
\noindent\emph{Source \texttt{CFT.pdf}, physical page 233.}\par\medskip
% BEGIN CFT_VII93_STATEMENT
\noindent\textsc{Lemma 9.3 (Key case of the Existence Theorem)}
\emph{Let $K$ be a number field containing a primitive $p$th root of $1$
($p$ prime). Then every open subgroup $V$ of $\boldsymbol{C}_K$ such
that $\boldsymbol{C}_K/V$ is a finite group killed by $p$ is a norm
group.}
% END CFT_VII93_STATEMENT

\section*{Original proof}
\noindent\emph{Source \texttt{CFT.pdf}, physical pages 233--234.}\par\medskip
% BEGIN CFT_VII93_PROOF
\noindent\textsc{Proof.} Let $S\supset S_\infty$ be a finite set of
primes of $K$ containing the infinite primes, those dividing $p$,
and enough primes so that $\mathbb{I}_K=K^\times\cdot\mathbb{I}_S$.
Let $L$ be the extension of $K$ corresponding by Kummer theory to the
group $U(S)\cdot K^{\times p}$, i.e., $L=K[U(S)^{\frac{1}{p}}]$, and let
\[
E=\prod_{v\in S}K_v^{\times p}\times\prod_{v\notin S}U_v.
\]
We shall prove that $K^\times\cdot E=K^\times\cdot
\operatorname{Nm}(\mathbb{I}_L)$ by verifying that
\begin{enumerate}
\renewcommand{\labelenumi}{(\alph{enumi})}
\item $E\subset\operatorname{Nm}(\mathbb{I}_L)$;
\item $(\mathbb{I}_K:K^\times\cdot E)=p^s=
(\mathbb{I}_K:K^\times\cdot\operatorname{Nm}_{L/K}(\mathbb{I}_L))$.
\end{enumerate}

For any prime $v$ of $K$ and prime $w$ of $L$ lying over it, the local
Artin map is an isomorphism
\[
K_v^\times/\operatorname{Nm}(L_w^\times)\longrightarrow
\operatorname{Gal}(L_w/K_v).
\]
Because $L/K$ is has exponent $p$,
$\operatorname{Nm}(L_w^\times)\supset K_v^{\times p}$.

For any prime $v\notin S$, $L$ is unramified over $K$ and $v$, and so
the norm map $U_w\longrightarrow U_v$ is onto.

On combining the statements in the last two paragraphs, we obtain (a).

From the Reciprocity Law,
\[
(\mathbb{I}_K:K^\times\cdot\operatorname{Nm}(\mathbb{I}_L))=[L:K],
\]
and from Kummer theory,
\[
[L:K]=(U(S)\cdot K^{\times p}:K^{\times p}).
\]
But
\[
U(S)\cdot K^{\times p}/K^{\times p}
\approx U(S)/U(S)\cap K^{\times p}.
\]
If $a^p\in U(S)$, then $a\in U(S)$, and so
$U(S)\cap K^{\times p}=U(S)^p$. Now, by the Dirichlet Unit Theorem
(ANT, 5.9),
\[
U(S)\approx U(S)_{\mathrm{torsion}}\oplus\mathbb{Z}^{s-1},
\]
and $U(S)_{\mathrm{torsion}}$ is the group of roots of $1$ in $K$,
which is a cyclic group whose order is divisible by $p$. Hence
$(U(S):U(S)^p)=p^s$.

On the other hand,
\[
(\mathbb{I}_K:K^\times\cdot E)
=(\mathbb{I}_S\cdot K^\times:E\cdot K^\times),
\]
which, by (6.5), equals
\[
(\mathbb{I}_S:E)/(\mathbb{I}_S\cap K^\times:E\cap K^\times).
\]
Therefore (see 6.8),
\[
(\mathbb{I}_S:E)=\prod_{v\in S}(K_v^\times:K_v^{\times p})
=\prod_{v\in S}\frac{p}{|p|_v}p=p^{2s}.
\]
Here, we have used that $K$ contains a primitive $p$th root of $1$ and
that $S$ contains all $v$ for which $|p|_v\ne1$, and so
$\prod_{v\in S}|p|_v=\prod_{\text{all }v}|p|_v=1$ by the product
formula. It follows that
$K^\times\cdot E=K^\times\cdot\operatorname{Nm}\mathbb{I}_L$.

Now let $V$ be an open subgroup of $\boldsymbol{C}_K$ such that
$\boldsymbol{C}_K/V$ is killed by $p$, and let $U$ be the inverse image
of $V$ in $\mathbb{I}_K$. Then $U$ is open in $\mathbb{I}_K$ and so
there is a finite set of primes $S$ such that
$U\supset\prod_{v\in S}1\times\prod_{v\notin S}U_v$.
Moreover, $\mathbb{I}_K/U$ has exponent $p$, and so
$U\supset\mathbb{I}_K^p$. Hence $U\supset E\cdot K^\times$, and because
$E\cdot K^\times/K^\times$ is a norm group, so also must be
$U/K^\times=V$. \hfill$\square$
% END CFT_VII93_PROOF

\section*{Original supporting excerpts}
\noindent\emph{Source \texttt{CFT.pdf}, physical pages 232--233.}\par\medskip
% BEGIN CFT_VII93_CONTEXT
\noindent\textsc{Lemma 9.1}
\emph{If $U$ is a norm group, and $V\supset U$, then $V$ also is a
norm group.}

\noindent\textsc{Proof.} Suppose $U=\operatorname{Nm}\boldsymbol{C}_L$.
According to the Reciprocity Law, the Artin map defines an isomorphism
\[
\phi:\boldsymbol{C}_K/U\longrightarrow\operatorname{Gal}(L/K).
\]
If $M$ is the fixed field of $\phi(V)$, then $\phi$ defines an
isomorphism
\[
\boldsymbol{C}_K/V\longrightarrow\operatorname{Gal}(M/K),
\]
but, according to the Reciprocity Law, the kernel of
$\phi:\boldsymbol{C}\longrightarrow\operatorname{Gal}(M/K)$ is
$\operatorname{Nm}_{M/K}\boldsymbol{C}_M$. \hfill$\square$

\medskip
\noindent\emph{Source \texttt{CFT.pdf}, physical pages 226--227.}\par\medskip
\noindent\textsc{Proposition 6.8}
\emph{Let $K$ be a local field of characteristic zero, and let
$\mu_n$ be the group of $n$th roots of $1$ in $K^\times$. Then}
\[
(K^\times:K^{\times n})=n\frac{|\mu_n|}{|n|}.
\]
\emph{If $K$ is nonarchimedean and $U$ is the group of units in $K$, then}
\[
(U:U^n)=\frac{|\mu_n|}{|n|}.
\]
\noindent\textsc{Proof.} If $K=\mathbb C$, the first equation becomes
$1=n\frac{n}{n^2}$, and if $K=\mathbb R$, it becomes
$1=n\frac{1}{n}$ when $n$ is odd and $2=n\frac{2}{n}$ when $n$ is even.
If $K$ is nonarchimedean, $K^\times\approx U\times\mathbb Z$, and so
\[
(K^\times:K^{\times n})=(U:U^n)(\mathbb Z:n\mathbb Z)=(U:U^n)n.
\]
Therefore, it remains to prove the second equation.

For an abelian group $M$, we write
\[
h_n(M)=(M:nM)/(M_n:1),\qquad M_n=\{x\in M\mid nx=0\}.
\]
Then $h_n(M)$ is the Herbrand quotient of $M$ regarded as a
$\mathbb Z/n\mathbb Z$-module with trivial action, and so we may apply
the results in II, \S2.

As we saw in the proof of (III, 2.4), the exponential map defines an
isomorphism from a subgroup of finite index in $\mathcal O_K$ onto a
subgroup of finite index in $U$. Therefore
\[
h_n(U)=h_n(\mathcal O_K)=(\mathcal O_K:n\mathcal O_K)
\overset{\mathrm{def}}{=}|n|^{-1}.
\]
Hence
\[
(U:U^n)=\frac{(U_n:1)}{|n|}.
\]
As $U_n=\mu_n$, this proves the second equation. \hfill$\square$

\medskip
\noindent\emph{Source \texttt{CFT.pdf}, physical page 235.}\par\medskip
\noindent\textsc{Theorem A.3} \emph{The map}
\[
L\longmapsto K^\times\cap L^{\times n}/K^{\times n}
\]
\emph{defines a one-to-one correspondence between the finite abelian
extensions of $K$ of exponent $n$ contained in some fixed algebraic
closure $\Omega$ of $K$ and the finite subgroups $B$ of
$K^\times/K^{\times n}$. The extension corresponding to $B$ is
$K[B^{\frac{1}{n}}]$, the smallest subfield of $\Omega$ containing $K$
and an $n$th root of each element of $B$. If $L\leftrightarrow B$,
then $[L:K]=(B:K^{\times n})$.}

\noindent\textsc{Proof.} See FT 5.29. \hfill$\square$
% END CFT_VII93_CONTEXT

\section*{Adjacent original result for the index comparison}
\noindent\emph{Source \texttt{CFT.pdf}, physical page 233. This excerpt
is separate from the selected proof of Lemma 9.3.}\par\medskip
% BEGIN CFT_VII93_PROP92
\noindent\textsc{Proposition 9.2}
\emph{Let $K$ be a number field containing a primitive $n$th root of $1$,
and let $S\supset S_\infty$ be a finite set of primes of $K$ containing
all those dividing $n$ and enough primes to generate the class group of
$K$. Any $a\in K^\times$ such that}
\begin{itemize}
\item[$\diamond$] \emph{$a$ is an $n$th power in $K_v$ for all $v\in S$;}
\item[$\diamond$] \emph{$a$ is a unit in $K_v$ for all $v\notin S$.}
\end{itemize}
\emph{is an $n$th power in $K$.}

\noindent\textsc{Proof.}\textsuperscript{7} Let $L=K[a^{1/n}]$---because
$\zeta_n\in K$, this is an abelian extension of $K$. For every prime
$v\in S$, $X^n-a$ splits completely in $K_v[X]$, and so $L_w=K_v$ for
all $w\mid v$. Hence the norm map $L_w^\times\longrightarrow K_v^\times$
is onto. On the other hand, $L$ is unramified over $K$ at any prime
$v\notin S$, and so the norm map
$\operatorname{Nm}_{L/K}:U_w\longrightarrow U_v$ is onto. Therefore,
$\operatorname{Nm}_{L/K}(\mathbb{I}_L)\supset\mathbb{I}_S$, and so
\[
K^\times\cdot\operatorname{Nm}_{L/K}(\mathbb{I}_L)
\supset K^\times\cdot\mathbb{I}_S=\mathbb{I}_K.
\]
The Reciprocity Law now shows that $L=K$, and so $a$ is an $n$th power
in $K$. \hfill$\square$

\medskip
\noindent{\small \textsuperscript{7}From Wojtek Wawr\'ow: Since $L=K$
in the proof is cyclic, it appears that the proof only uses the First
Inequality (specifically Lemma 4.5), not the full Reciprocity Law.
Regardless, it appears that this Proposition is not used anywhere.}
% END CFT_VII93_PROP92
\end{document}
